% ============================================================================
% Fichier  : partie1/P01_C02_risque_decision.tex
% Titre    : Du Risque à la Décision
% Auteur   : MINKA MI NGUIDJOÏ Thierry Emmanuel
% ============================================================================
% Sources : ISO/IEC 27001:2022 ; ISO 31000:2018. Mairie d'Eséka : cas fictif.
% ============================================================================

\chapter{Du Risque à la Décision}
\label{chap:enjeux}
\chaptermark{Du risque à la décision}

\epigraph{%
    « Sécuriser, c'est décider~: quel risque accepter, lequel réduire,\\
    et avec quelles ressources. »%
}{Principe du cours}

\niveauChapitre{Fondamental}{Le Chapitre~\ref{chap:cia}.}

\begin{boxobjectifs}
\begin{itemize}
    \item Relier la sécurité aux enjeux économiques, réglementaires et
          stratégiques de l'organisation.
    \item Identifier les acteurs de la sécurité et leurs responsabilités.
    \item Définir le patrimoine informationnel à protéger.
    \item Comprendre que la sécurité est une décision fondée sur le risque.
    \item Situer la mairie d'Eséka dans son écosystème (ANTIC, parties
          prenantes).
\end{itemize}
\end{boxobjectifs}

\bigskip

% ============================================================================
\section*{Question directrice : au nom de quoi la mairie investit-elle dans la sécurité ?}
% ============================================================================

Dépenser pour la sécurité, c'est renoncer à dépenser ailleurs. Ce chapitre
montre que cette dépense se justifie par des enjeux, s'organise autour
d'acteurs et se décide en fonction du risque.

\medskip
\textbf{Plan du chapitre.}
\begin{enumerate}
    \item Les enjeux (section~\ref{sec:c02-enjeux}).
    \item Les acteurs (section~\ref{sec:c02-acteurs}).
    \item Le patrimoine informationnel (section~\ref{sec:c02-patrimoine}).
    \item La sécurité, une décision (section~\ref{sec:c02-decision}).
    \item L'écosystème de la mairie (section~\ref{sec:c02-ecosysteme}).
\end{enumerate}

\separateur

% ============================================================================
\section{Les Enjeux}
\label{sec:c02-enjeux}
% ============================================================================

\begin{tableaucours}{L{3.4cm}L{9.0cm}}{%
    \tableauHeader{Enjeu} &
    \tableauHeader{Pour la mairie d'Eséka}
}
Économique & Service interrompu, recettes de taxes retardées, remédiation coûteuse \\
Réglementaire & Loi 2024/017 sur les données personnelles, obligations de l'ANTIC \\
Stratégique & Confiance des usagers, image de la commune, continuité du service public \\
\end{tableaucours}
\vspace{-0.8em}
\begin{center}{\small\itshape Tableau~2.1. Les enjeux de la sécurité pour la mairie}\end{center}

Pour une collectivité, l'enjeu premier n'est pas le chiffre d'affaires mais la
\termecle{continuité du service public} et la confiance des administrés. Un
état civil à l'arrêt, c'est un citoyen qui ne peut ni se marier, ni inscrire
un enfant, ni prouver son identité.

% ============================================================================
\section{Les Acteurs}
\label{sec:c02-acteurs}
% ============================================================================

\begin{tableaucours}{L{3.6cm}L{8.8cm}}{%
    \tableauHeader{Acteur} &
    \tableauHeader{Responsabilité}
}
Le maire (direction) & Impulse, valide la politique, alloue les ressources, assume \\
RSSI & Pilote la sécurité, apprécie les risques, propose les mesures \\
DSI & Met en \oe uvre les solutions techniques \\
DPO / conseil & Veille à la conformité des traitements de données personnelles \\
Agents (utilisateurs) & Respectent les règles, signalent les incidents \\
ANTIC & Autorité nationale : appui, contrôle, réponse aux incidents \\
\end{tableaucours}
\vspace{-0.8em}
\begin{center}{\small\itshape Tableau~2.2. Les acteurs et leurs responsabilités}\end{center}

La sécurité n'est l'affaire de personne si elle est l'affaire de tous sans
responsabilité nommée. La première décision du maire, nommer une \sigle{RSSI},
crée le point où la responsabilité se concentre.

\begin{boxmetier}{Direction générale}
\textbf{Le rôle que nul ne peut déléguer.} Le maire peut déléguer la mise en
\oe uvre, jamais la responsabilité finale. C'est lui qui arbitre entre le
risque et le coût, parce que lui seul engage la commune.
\end{boxmetier}

% ============================================================================
\section{Le Patrimoine Informationnel}
\label{sec:c02-patrimoine}
% ============================================================================

\begin{boxdef}
\termecle{Patrimoine informationnel}~: ensemble des informations de valeur
d'une organisation et des moyens qui les portent. Le protéger suppose d'abord
de le connaître.
\end{boxdef}

\begin{boxobjet}{Le patrimoine de la mairie}
\textbf{Constatation.} Avant la démarche, la mairie n'a aucun inventaire de
ses informations sensibles.

\medskip
\textbf{Analyse.} On ne protège bien que ce que l'on a recensé. L'opération
\textsc{Nyong} commencera par un inventaire~: état civil, fichier des
contribuables, délibérations du conseil, messagerie, site web. Sans cet
inventaire, l'appréciation des risques n'a pas d'objet.
\end{boxobjet}

% ============================================================================
\section{La Sécurité, une Décision}
\label{sec:c02-decision}
% ============================================================================

La sécurité absolue n'existe pas et coûterait trop cher. Sécuriser, c'est
donc \termecle{décider}~: pour chaque risque, choisir de le réduire, de le
transférer, de l'éviter ou de l'accepter (Chapitre~\ref{chap:traitement}).
Cette décision appartient à la direction, sur proposition du \sigle{RSSI} et au
vu des risques appréciés.

\begin{boxinfo}
\textbf{Le fondement du SMSI.} Tout le cours découle de cette idée~: la
sécurité se \emph{décide} sur la base du \emph{risque}. C'est pourquoi
l'appréciation des risques (partie IV) est le cœur de la norme
\norme{ISO/IEC 27001}.
\end{boxinfo}

% ============================================================================
\section{L'Écosystème de la Mairie}
\label{sec:c02-ecosysteme}
% ============================================================================

La mairie n'est pas seule. Elle dépend de prestataires (hébergeur du site,
éditeur du logiciel d'état civil), relève d'autorités (ANTIC, tutelle) et sert
des usagers. Ces \termecle{parties intéressées} ont des attentes de sécurité
qui s'imposent au \sigle{SMSI} et seront recensées à la clause~4 de la norme
(Chapitre~\ref{chap:clauses46}).

\begin{boxjuris}
La loi camerounaise n°~2024/017 du 23~décembre~2024 sur la protection des
données à caractère personnel impose au responsable de traitement (ici, la
mairie) de garantir la sécurité des données qu'il détient. La sécurité n'est
donc pas seulement prudente~: elle est une obligation légale.
\end{boxjuris}

\separateur

\begin{boxresume}
\textbf{Ce qu'il faut retenir de ce chapitre~:}
\begin{itemize}
    \item La sécurité répond à des \textbf{enjeux} économiques,
          réglementaires et stratégiques~; pour une mairie, la continuité du
          service public prime.
    \item Chaque \textbf{acteur} a une responsabilité nommée~; la direction
          assume, le \sigle{RSSI} pilote.
    \item On protège un \textbf{patrimoine informationnel} qu'il faut d'abord
          inventorier.
    \item Sécuriser est une \textbf{décision} fondée sur le risque.
    \item Les \textbf{parties intéressées} imposent des attentes au
          \sigle{SMSI}.
\end{itemize}
\end{boxresume}

\begin{boxinfo}
\textbf{Points de contrôle de l'élève.} Je sais~: citer les trois types
d'enjeux~; nommer les acteurs et leurs rôles~; définir le patrimoine
informationnel~; expliquer pourquoi la sécurité est une décision~; citer une
obligation légale.
\end{boxinfo}

\bigskip

\begin{boxexo}[title={\ding{45}\quad Exercice~2.1. Trois enjeux \niveaubase}]
Donnez, pour la mairie, un exemple concret d'enjeu économique, réglementaire
et stratégique.
\end{boxexo}

\begin{boxexo}[title={\ding{45}\quad Exercice~2.2. Qui fait quoi \niveaubase}]
Reliez chaque acteur (maire, RSSI, DSI, agent) à sa responsabilité principale.
\end{boxexo}

\begin{boxexo}[title={\ding{45}\quad Exercice~2.3. Inventaire \niveaubase}]
Listez cinq informations de valeur que détient une mairie.
\end{boxexo}

\begin{boxexo}[title={\ding{45}\quad Exercice~2.4. Responsabilité finale \niveaumoyen}]
Pourquoi le maire ne peut-il pas déléguer la responsabilité finale de la
sécurité~?
\end{boxexo}

\begin{boxexo}[title={\ding{45}\quad Exercice~2.5. Décider \niveaumoyen}]
« La mairie n'a qu'à tout sécuriser au maximum. » Montrez les limites de cette
phrase.
\end{boxexo}

\begin{boxexo}[title={\ding{45}\quad Exercice~2.6. Parties intéressées \niveaumoyen}]
Identifiez quatre parties intéressées de la mairie et une attente de sécurité
pour chacune.
\end{boxexo}

\begin{boxexo}[title={\ding{45}\quad Exercice~2.7. Continuité publique \niveaumoyen}]
Pourquoi la continuité du service public pèse-t-elle plus, pour une mairie,
que pour une entreprise commerciale~?
\end{boxexo}

\begin{boxexo}[title={\ding{45}\quad Exercice~2.8. Obligation légale \niveaumoyen}]
En quoi la loi 2024/017 transforme-t-elle la sécurité des données de la mairie
en obligation~?
\end{boxexo}

\begin{boxexo}[title={\ding{45}\quad Exercice~2.9. Cartographie des acteurs \niveauexpert}]
Dressez l'organigramme de sécurité que la mairie d'Eséka devrait adopter, avec
les responsabilités de chacun.
\end{boxexo}

\begin{boxexo}[title={\ding{45}\quad Exercice~2.10. Argumentaire budgétaire \niveauexpert}]
Rédigez l'argumentaire par lequel la RSSI justifie, devant le conseil, un
budget sécurité, en s'appuyant sur les enjeux et le risque.
\end{boxexo}

\bigskip

\begin{boxinfo}
\textbf{Pour aller plus loin.}
\begin{itemize}
    \item ISO/IEC 27001~\cite{iso27001}~; ISO 31000~\cite{iso31000}.
\end{itemize}
\end{boxinfo}

\begin{boxinfo}
\textbf{Transition.} Décider une fois ne suffit pas~: il faut un cycle qui
entretienne la décision. Le Chapitre~\ref{chap:pdca} présente le \sigle{PDCA}.
\end{boxinfo}
